\section*{Evaluasi}
\addcontentsline{toc}{section}{Evaluasi}

\begin{enumerate}[label=\arabic*.,leftmargin=*]
  \item Definisikan tata bahasa \(G = (V,T,P,S)\) dan jelaskan arti
        setiap komponen.
  \item Diberikan \(L = \{w \in \{a,b\}^{*} \mid |w| \text{ genap}\}\).
        Tuliskan tata bahasa regular untuk \(L\).
  \item Jelaskan mengapa \(L = \{a^n b^n \mid n \ge 0\}\) tidak dapat
        dikenali DFA, lalu klasifikasikan bahasa tersebut pada
        hierarki Chomsky.
  \item Buat derivasi leftmost, derivasi rightmost, dan pohon urai
        untuk string \(a + b * c\) pada tata bahasa aritmetika.
  \item Jelaskan perbedaan daya ekspresif otomata hingga, otomata
        pushdown, dan mesin Turing.
  \item Mengapa analisis semantik tidak dapat dilakukan hanya dengan
        tata bahasa bebas konteks? Berikan satu contoh.
\end{enumerate}

\subsection*{Kunci dan Rubrik Singkat}
\begin{itemize}[leftmargin=*]
  \item Soal 1--2: ketepatan definisi dan notasi (70\%), kelengkapan
        (30\%).
  \item Soal 3: ketepatan penjelasan (60\%), klasifikasi (40\%).
  \item Soal 4: kebenaran derivasi (50\%), kebenaran pohon (50\%).
  \item Soal 5--6: ketepatan konsep (60\%), contoh (40\%).
\end{itemize}
